\documentclass[11pt,a4paper]{article}

\usepackage[utf8]{inputenc}
\usepackage[T1]{fontenc}
\usepackage{lmodern}
\usepackage[margin=1in]{geometry}
\usepackage{hyperref}
\usepackage{xcolor}
\usepackage{listings}
\usepackage{enumitem}
\usepackage{graphicx}
\usepackage{titlesec}
\usepackage{tcolorbox}
\usepackage{array}
\usepackage{booktabs}
\usepackage{longtable}

% ---------- Hyperref ----------
\hypersetup{
    colorlinks=true,
    linkcolor=blue!60!black,
    urlcolor=blue!60!black,
    citecolor=blue!60!black,
    pdftitle={Deploying ML Models on AWS Lambda + ECR},
    pdfauthor={Prabal}
}

% ---------- Listings ----------
\definecolor{codebg}{RGB}{245,245,245}
\definecolor{codekw}{RGB}{0,0,180}
\definecolor{codestr}{RGB}{163,21,21}
\definecolor{codecom}{RGB}{0,128,0}

\lstdefinestyle{pycode}{
    backgroundcolor=\color{codebg},
    basicstyle=\ttfamily\footnotesize,
    keywordstyle=\color{codekw}\bfseries,
    stringstyle=\color{codestr},
    commentstyle=\color{codecom}\itshape,
    language=Python,
    breaklines=true,
    showstringspaces=false,
    numbers=left,
    numberstyle=\tiny\color{gray},
    frame=single,
    rulecolor=\color{gray!30},
    tabsize=2
}

\lstdefinestyle{bashcode}{
    backgroundcolor=\color{codebg},
    basicstyle=\ttfamily\footnotesize,
    keywordstyle=\color{codekw}\bfseries,
    stringstyle=\color{codestr},
    commentstyle=\color{codecom}\itshape,
    language=bash,
    breaklines=true,
    showstringspaces=false,
    frame=single,
    rulecolor=\color{gray!30},
    tabsize=2
}

\lstdefinelanguage{docker}{
    keywords={FROM, RUN, CMD, ENTRYPOINT, WORKDIR, COPY, ADD, ENV, ARG, EXPOSE, VOLUME, USER, LABEL, HEALTHCHECK, SHELL},
    keywordstyle=\color{codekw}\bfseries,
    sensitive=true,
    comment=[l]{\#},
    morestring=[b]",
}

\lstdefinestyle{dockercode}{
    backgroundcolor=\color{codebg},
    basicstyle=\ttfamily\footnotesize,
    keywordstyle=\color{codekw}\bfseries,
    stringstyle=\color{codestr},
    commentstyle=\color{codecom}\itshape,
    language=docker,
    breaklines=true,
    showstringspaces=false,
    frame=single,
    rulecolor=\color{gray!30},
    tabsize=2
}

\lstdefinelanguage{json}{
    string=[s]{"}{"},
    stringstyle=\color{codestr},
    comment=[l]{//},
    commentstyle=\color{codecom},
}

\lstdefinestyle{jsoncode}{
    backgroundcolor=\color{codebg},
    basicstyle=\ttfamily\footnotesize,
    language=json,
    breaklines=true,
    showstringspaces=false,
    frame=single,
    rulecolor=\color{gray!30},
    tabsize=2
}

% ---------- Note boxes ----------
\newtcolorbox{notebox}{
    colback=blue!5!white, colframe=blue!50!black,
    title=Note, fonttitle=\bfseries, sharp corners, boxrule=0.5pt
}
\newtcolorbox{warnbox}{
    colback=orange!5!white, colframe=orange!70!black,
    title=Important, fonttitle=\bfseries, sharp corners, boxrule=0.5pt
}
\newtcolorbox{tipbox}{
    colback=green!5!white, colframe=green!50!black,
    title=Tip, fonttitle=\bfseries, sharp corners, boxrule=0.5pt
}

\titleformat{\section}{\Large\bfseries\color{blue!50!black}}{\thesection}{1em}{}
\titleformat{\subsection}{\large\bfseries\color{blue!40!black}}{\thesubsection}{1em}{}

\title{\textbf{Deploying ML Models on AWS Lambda with ECR}\\
\large A Step-by-Step Serverless Deployment Guide}
\author{Prabal}
\date{\today}

\begin{document}
\maketitle
\tableofcontents
\newpage

% =========================================================
\section{Introduction}
% =========================================================

AWS Lambda is a serverless compute service that runs code on demand without
provisioning or managing servers. When combined with Amazon Elastic Container
Registry (ECR), Lambda can execute \textbf{container images up to 10 GB},
which is large enough to host many small-to-medium ML models along with their
inference runtime (PyTorch, ONNX, scikit-learn, HuggingFace Transformers).

This deployment pattern --- often called \emph{``Lambda + ECR''} --- is the
most cost-effective way to serve ML models that receive \textbf{sporadic or
unpredictable traffic}. You pay only for the milliseconds during which your
function is actually executing; idle endpoints cost nothing.

\begin{notebox}
\textbf{When to prefer Lambda + ECR over SageMaker:}
\begin{itemize}[leftmargin=*,nosep]
    \item Model fits in $\leq$ 10 GB image and $\leq$ 10 GB memory.
    \item Traffic is spiky, infrequent, or unpredictable (hours of idle per day).
    \item Inference latency tolerance is moderate (cold starts can be seconds).
    \item No GPU required (Lambda is CPU-only).
    \item You want tight integration with API Gateway, S3 events, EventBridge,
          SQS, Step Functions, etc.
\end{itemize}
\end{notebox}

% =========================================================
\section{Architecture Overview}
% =========================================================

The typical topology:

\begin{center}
\texttt{Client $\rightarrow$ API Gateway $\rightarrow$ Lambda (container) $\rightarrow$ Response}
\end{center}

with the container image pulled from ECR at cold start, and model weights
optionally cached from S3 or baked directly into the image.

\textbf{Key components:}
\begin{itemize}[leftmargin=*,nosep]
    \item \textbf{ECR} --- private Docker registry that stores the inference image.
    \item \textbf{Lambda} --- executes the container on demand; reuses warm
          containers across invocations.
    \item \textbf{API Gateway} / \textbf{Function URL} --- HTTPS entry point.
    \item \textbf{IAM role} --- Lambda execution role granting permission to
          pull from ECR and read model artifacts from S3 (if externalized).
    \item \textbf{CloudWatch Logs \& Metrics} --- observability.
\end{itemize}

% =========================================================
\section{Lambda Constraints (Read First)}
% =========================================================

Before committing to this path, confirm your model fits within Lambda's limits:

\begin{longtable}{p{5cm} p{9.5cm}}
\toprule
\textbf{Resource} & \textbf{Hard limit} \\
\midrule
Container image size & 10 GB (uncompressed, as stored in ECR) \\
\addlinespace
Memory & 128 MB -- 10 240 MB (10 GB) \\
\addlinespace
vCPU & Scales linearly with memory (6 vCPUs at 10 GB) \\
\addlinespace
Execution timeout & 15 minutes (900 s) \\
\addlinespace
\texttt{/tmp} scratch space & 512 MB -- 10 GB (configurable) \\
\addlinespace
Request / response payload & 6 MB sync, 256 KB async (API Gateway: 10 MB) \\
\addlinespace
Concurrent executions & 1 000 per region by default (soft limit) \\
\addlinespace
GPU & Not supported \\
\bottomrule
\end{longtable}

\begin{warnbox}
If your model weights + runtime exceed \textbf{10 GB}, or if you need GPU
inference, Lambda is the wrong tool --- use SageMaker, ECS/Fargate, or EC2
instead.
\end{warnbox}

% =========================================================
\section{Prerequisites}
% =========================================================

\begin{enumerate}[leftmargin=*]
    \item AWS account, AWS CLI configured (\texttt{aws configure}).
    \item Docker installed and running locally.
    \item Python 3.9+.
    \item A trained model artifact (e.g., \texttt{model.pkl},
          \texttt{model.onnx}, \texttt{model.pt}, or a HuggingFace checkpoint
          folder).
    \item IAM permissions to create ECR repos, Lambda functions, and IAM roles.
\end{enumerate}

% =========================================================
\section{Step-by-Step Deployment}
% =========================================================

% --------------------------------------------
\subsection{Step 1 --- Project Layout}
% --------------------------------------------

A minimal layout:

\begin{lstlisting}[style=bashcode]
lambda-ml/
  app.py               # Lambda handler
  requirements.txt     # Python dependencies
  model/               # model weights (baked into image)
    model.pkl
  Dockerfile
\end{lstlisting}

% --------------------------------------------
\subsection{Step 2 --- Write the Lambda Handler}
% --------------------------------------------

The handler loads the model \emph{once at module scope} so it survives across
warm invocations. Loading inside \texttt{lambda\_handler} would pay the cost
on every request.

\begin{lstlisting}[style=pycode]
import json
import pickle
import os
import numpy as np

# ---- Cold-start initialization (runs once per container) ----
MODEL_PATH = os.environ.get("MODEL_PATH", "/var/task/model/model.pkl")
with open(MODEL_PATH, "rb") as f:
    MODEL = pickle.load(f)

def lambda_handler(event, context):
    try:
        # API Gateway proxy integration wraps the payload in "body"
        body = event.get("body", event)
        if isinstance(body, str):
            body = json.loads(body)

        features = np.array(body["inputs"], dtype=np.float32)
        preds = MODEL.predict(features).tolist()

        return {
            "statusCode": 200,
            "headers": {"Content-Type": "application/json"},
            "body": json.dumps({"predictions": preds}),
        }
    except Exception as e:
        return {
            "statusCode": 500,
            "body": json.dumps({"error": str(e)}),
        }
\end{lstlisting}

\begin{tipbox}
The model is loaded at import time so it is cached across invocations on the
same warm container. AWS keeps warm containers alive for $\sim$5--15 minutes of
inactivity, and longer under sustained traffic.
\end{tipbox}

% --------------------------------------------
\subsection{Step 3 --- Write the Dockerfile}
% --------------------------------------------

Use the official AWS Lambda Python base image --- it ships the Lambda Runtime
Interface Client (RIC) pre-installed, so your container already speaks the
Lambda invocation protocol.

\begin{lstlisting}[style=dockercode]
# Base image maintained by AWS (Python 3.11)
FROM public.ecr.aws/lambda/python:3.11

# Install Python dependencies into the Lambda task root
COPY requirements.txt .
RUN pip install --no-cache-dir -r requirements.txt \
    --target "${LAMBDA_TASK_ROOT}"

# Copy code and model weights
COPY app.py        ${LAMBDA_TASK_ROOT}/
COPY model/        ${LAMBDA_TASK_ROOT}/model/

# Tell Lambda which function to invoke:  file.function
CMD ["app.lambda_handler"]
\end{lstlisting}

\textbf{Example \texttt{requirements.txt}:}
\begin{lstlisting}[style=bashcode]
numpy==1.26.4
scikit-learn==1.4.2
\end{lstlisting}

\begin{warnbox}
The base image is \texttt{linux/amd64}. If you are on an Apple Silicon Mac,
build with \texttt{--platform linux/amd64} or the image will fail to start in
Lambda with cryptic \texttt{exec format error}.
\end{warnbox}

% --------------------------------------------
\subsection{Step 4 --- Create an ECR Repository}
% --------------------------------------------

\begin{lstlisting}[style=bashcode]
export AWS_REGION=us-east-1
export ACCOUNT_ID=$(aws sts get-caller-identity --query Account --output text)
export REPO_NAME=lambda-ml-inference

aws ecr create-repository \
    --repository-name $REPO_NAME \
    --region $AWS_REGION \
    --image-scanning-configuration scanOnPush=true
\end{lstlisting}

% --------------------------------------------
\subsection{Step 5 --- Build, Tag, Push the Image}
% --------------------------------------------

\begin{lstlisting}[style=bashcode]
# 1. Authenticate Docker to ECR
aws ecr get-login-password --region $AWS_REGION | \
    docker login --username AWS \
    --password-stdin $ACCOUNT_ID.dkr.ecr.$AWS_REGION.amazonaws.com

# 2. Build (force linux/amd64 for Lambda)
docker build --platform linux/amd64 -t $REPO_NAME:v1 .

# 3. Tag for ECR
docker tag $REPO_NAME:v1 \
    $ACCOUNT_ID.dkr.ecr.$AWS_REGION.amazonaws.com/$REPO_NAME:v1

# 4. Push
docker push $ACCOUNT_ID.dkr.ecr.$AWS_REGION.amazonaws.com/$REPO_NAME:v1
\end{lstlisting}

\begin{tipbox}
Use immutable, version-tagged images (\texttt{v1}, \texttt{v2}, git SHA) rather
than \texttt{latest}. Lambda pins the image digest at function creation, so
overwriting \texttt{latest} does not actually roll out a new version --- and
worse, makes rollbacks ambiguous.
\end{tipbox}

% --------------------------------------------
\subsection{Step 6 --- Create the Lambda Execution Role}
% --------------------------------------------

\textbf{Trust policy} (\texttt{trust-policy.json}):

\begin{lstlisting}[style=jsoncode]
{
  "Version": "2012-10-17",
  "Statement": [{
    "Effect": "Allow",
    "Principal": {"Service": "lambda.amazonaws.com"},
    "Action": "sts:AssumeRole"
  }]
}
\end{lstlisting}

Create the role and attach the basic execution policy:

\begin{lstlisting}[style=bashcode]
aws iam create-role \
    --role-name lambda-ml-role \
    --assume-role-policy-document file://trust-policy.json

aws iam attach-role-policy \
    --role-name lambda-ml-role \
    --policy-arn arn:aws:iam::aws:policy/service-role/AWSLambdaBasicExecutionRole
\end{lstlisting}

Attach extra permissions if the function needs to read weights from S3:

\begin{lstlisting}[style=bashcode]
aws iam attach-role-policy \
    --role-name lambda-ml-role \
    --policy-arn arn:aws:iam::aws:policy/AmazonS3ReadOnlyAccess
\end{lstlisting}

% --------------------------------------------
\subsection{Step 7 --- Create the Lambda Function}
% --------------------------------------------

\begin{lstlisting}[style=bashcode]
aws lambda create-function \
    --function-name ml-inference \
    --package-type Image \
    --code ImageUri=$ACCOUNT_ID.dkr.ecr.$AWS_REGION.amazonaws.com/$REPO_NAME:v1 \
    --role arn:aws:iam::$ACCOUNT_ID:role/lambda-ml-role \
    --timeout 60 \
    --memory-size 3008 \
    --ephemeral-storage Size=2048 \
    --architectures x86_64 \
    --region $AWS_REGION
\end{lstlisting}

\begin{notebox}
\textbf{Memory sizing is critical.} Lambda scales vCPU linearly with memory ---
more memory often makes the function \emph{both faster and cheaper} per
invocation because billed duration drops. For ML inference, start at
\textbf{3 008 MB (2 vCPU)} and benchmark up/down from there.
\end{notebox}

% --------------------------------------------
\subsection{Step 8 --- Expose via HTTP}
% --------------------------------------------

Two simple options:

\paragraph{Option A --- Function URL (quickest).}

\begin{lstlisting}[style=bashcode]
aws lambda create-function-url-config \
    --function-name ml-inference \
    --auth-type NONE    # or AWS_IAM for signed requests

aws lambda add-permission \
    --function-name ml-inference \
    --statement-id FunctionURLAllowPublicAccess \
    --action lambda:InvokeFunctionUrl \
    --principal "*" \
    --function-url-auth-type NONE
\end{lstlisting}

Returns a URL like:
\texttt{https://<random>.lambda-url.us-east-1.on.aws/}

\paragraph{Option B --- API Gateway (production-grade).}
Provides throttling, custom domains, WAF, usage plans, and request/response
transformation. Attach an HTTP API with a Lambda proxy integration pointing to
the function ARN.

% --------------------------------------------
\subsection{Step 9 --- Invoke}
% --------------------------------------------

\begin{lstlisting}[style=bashcode]
curl -X POST https://<function-url>/ \
     -H "Content-Type: application/json" \
     -d '{"inputs": [[0.1, 0.4, 0.9]]}'
\end{lstlisting}

Or from Python:

\begin{lstlisting}[style=pycode]
import boto3, json
lam = boto3.client("lambda", region_name="us-east-1")
resp = lam.invoke(
    FunctionName="ml-inference",
    Payload=json.dumps({"inputs": [[0.1, 0.4, 0.9]]}),
)
print(json.loads(resp["Payload"].read()))
\end{lstlisting}

% --------------------------------------------
\subsection{Step 10 --- Update and Roll Back}
% --------------------------------------------

Push a new tag, then point Lambda at the new digest:

\begin{lstlisting}[style=bashcode]
docker build --platform linux/amd64 -t $REPO_NAME:v2 .
docker tag  $REPO_NAME:v2 \
    $ACCOUNT_ID.dkr.ecr.$AWS_REGION.amazonaws.com/$REPO_NAME:v2
docker push $ACCOUNT_ID.dkr.ecr.$AWS_REGION.amazonaws.com/$REPO_NAME:v2

aws lambda update-function-code \
    --function-name ml-inference \
    --image-uri $ACCOUNT_ID.dkr.ecr.$AWS_REGION.amazonaws.com/$REPO_NAME:v2
\end{lstlisting}

Use \textbf{versions and aliases} plus \textbf{CodeDeploy} for canary / linear
rollouts with automatic rollback on CloudWatch alarm.

% =========================================================
\section{Dealing with Cold Starts}
% =========================================================

Cold start = the time Lambda spends pulling the image, initializing the
container, and running module-level code (including model load). For ML images
this is often \textbf{2--15 seconds}, which can dominate tail latency.

\textbf{Mitigations, in order of impact:}

\begin{enumerate}[leftmargin=*]
    \item \textbf{Provisioned Concurrency} --- reserves a pool of always-warm
          containers; eliminates cold starts entirely within that pool.
          Increases cost (billed hourly for reserved capacity).
    \item \textbf{SnapStart for Python / containers} --- caches initialized
          state snapshots; dramatically reduces cold-start for supported
          runtimes.
    \item \textbf{Smaller images} --- prune dev dependencies; use slim base
          images; quantize or compress model weights.
    \item \textbf{Lazy imports} --- only \texttt{import torch} etc. inside
          the functions that need them if the handler has multiple code paths.
    \item \textbf{Right-size memory upward} --- more memory = more CPU = faster
          init. Counter-intuitively often \emph{cheaper} because duration drops.
    \item \textbf{Pre-warming pings} --- EventBridge rule hitting the function
          every few minutes. Cheap and cheerful but not guaranteed.
\end{enumerate}

\begin{tipbox}
For customer-facing synchronous APIs on spiky traffic, \textbf{Provisioned
Concurrency} is usually the right answer. Reserve 1--2 instances; it is
inexpensive compared to a SageMaker endpoint running 24/7.
\end{tipbox}

% =========================================================
\section{Externalizing Model Weights}
% =========================================================

Baking weights into the image is simplest but rebuilds on every model update.
For faster iteration or for models that approach the 10 GB image limit, store
weights in S3 and fetch them at cold start into \texttt{/tmp}:

\begin{lstlisting}[style=pycode]
import os, boto3, pickle

BUCKET = os.environ["MODEL_BUCKET"]
KEY    = os.environ["MODEL_KEY"]
LOCAL  = "/tmp/model.pkl"

if not os.path.exists(LOCAL):
    boto3.client("s3").download_file(BUCKET, KEY, LOCAL)

with open(LOCAL, "rb") as f:
    MODEL = pickle.load(f)
\end{lstlisting}

\begin{notebox}
\texttt{/tmp} persists across warm invocations but \emph{not} across cold
starts. Each new container will re-download. For multi-GB weights, this can
add seconds to cold start --- weigh the tradeoff against image-based loading.
For shared model storage across many Lambda containers, consider mounting
Amazon EFS.
\end{notebox}

% =========================================================
\section{Framework-Specific Notes}
% =========================================================

\subsection{PyTorch}
\begin{itemize}[leftmargin=*,nosep]
    \item Install the \emph{CPU-only} wheel: \texttt{torch --index-url
          https://download.pytorch.org/whl/cpu} --- avoids CUDA bloat and
          keeps the image under the 10 GB limit.
    \item Call \texttt{model.eval()} and wrap inference in
          \texttt{torch.inference\_mode()}.
    \item Consider exporting to TorchScript or ONNX for faster load and
          inference.
\end{itemize}

\subsection{HuggingFace Transformers}
\begin{itemize}[leftmargin=*,nosep]
    \item Set \texttt{TRANSFORMERS\_CACHE=/tmp} to avoid writes into the
          read-only image layer.
    \item Pre-download weights at build time with a \texttt{RUN python -c
          "..."} step rather than fetching from the Hub at cold start.
    \item For large models (7B+) prefer quantized GGUF / int4 variants or
          skip Lambda entirely.
\end{itemize}

\subsection{ONNX Runtime}
\begin{itemize}[leftmargin=*,nosep]
    \item Excellent fit: small runtime, fast CPU inference, low cold start.
    \item Use \texttt{onnxruntime} (not \texttt{onnxruntime-gpu}) for Lambda.
\end{itemize}

\subsection{scikit-learn / XGBoost / LightGBM}
\begin{itemize}[leftmargin=*,nosep]
    \item Near-ideal fit --- small dependencies, fast cold start, cheap.
    \item Pickle version pinning matters: load with the same library version
          that pickled the model.
\end{itemize}

% =========================================================
\section{Observability}
% =========================================================

\begin{itemize}[leftmargin=*]
    \item \textbf{CloudWatch Logs} --- all \texttt{stdout}/\texttt{stderr}
          lands in \texttt{/aws/lambda/<function-name>}. Use structured JSON
          logs for easy querying with CloudWatch Logs Insights.
    \item \textbf{CloudWatch Metrics} --- \texttt{Invocations}, \texttt{Errors},
          \texttt{Throttles}, \texttt{Duration}, \texttt{ConcurrentExecutions},
          \texttt{IteratorAge} (streams).
    \item \textbf{AWS X-Ray} --- enable active tracing for per-segment latency
          breakdown (init vs. invoke vs. downstream calls).
    \item \textbf{Lambda Insights} --- enhanced metrics (CPU, memory,
          network, disk) via a bundled extension layer.
\end{itemize}

% =========================================================
\section{Cost Model}
% =========================================================

Lambda pricing has two components per invocation:

\begin{enumerate}[leftmargin=*,nosep]
    \item \textbf{Requests} --- \$0.20 per 1 M requests.
    \item \textbf{Duration} --- billed per ms, proportional to memory
          allocated (GB-seconds).
\end{enumerate}

\paragraph{Back-of-envelope example.} A function at 3 008 MB averaging 500 ms
per invocation, 1 M invocations per month:
\[
\text{Duration cost} \approx 1{,}000{,}000 \times 0.5\,\text{s}
\times \frac{3008}{1024}\,\text{GB} \times \$0.0000166667/\text{GB-s}
\approx \$24.50
\]
plus \$0.20 requests = $\sim$\$24.70/month. Compared to a \texttt{ml.m5.xlarge}
SageMaker endpoint at $\sim$\$140/month running 24/7, this is $\sim$6$\times$
cheaper for the same invocation count --- but SageMaker wins once traffic is
high enough to keep the endpoint continuously busy.

\begin{tipbox}
The crossover point where SageMaker becomes cheaper is typically around
\textbf{5--10 million invocations per month} for CPU workloads. Below that,
Lambda wins; above that, a right-sized persistent endpoint wins.
\end{tipbox}

% =========================================================
\section{Production Best Practices}
% =========================================================

\begin{enumerate}[leftmargin=*]
    \item \textbf{Infrastructure as Code} --- Define everything (ECR repo,
          Lambda, IAM role, API Gateway) in AWS SAM, CDK, or Terraform.
    \item \textbf{CI/CD} --- GitHub Actions / CodePipeline: build image
          $\rightarrow$ push to ECR $\rightarrow$ update function $\rightarrow$
          shift alias with CodeDeploy canary.
    \item \textbf{Versions \& aliases} --- Publish immutable versions; point
          \texttt{prod} and \texttt{staging} aliases at specific versions for
          safe rollouts.
    \item \textbf{Idempotency} --- For async / SQS-triggered inference,
          de-duplicate on a request ID; Lambda may retry on failure.
    \item \textbf{Reserved / Provisioned Concurrency} --- Protect critical
          functions from noisy-neighbor throttling; eliminate cold starts.
    \item \textbf{VPC only when needed} --- VPC-attached Lambdas add ENI
          cold-start overhead. Attach only if you must reach RDS / ElastiCache
          / private subnets. Use VPC endpoints to avoid NAT Gateway charges.
    \item \textbf{Secrets} --- Use Secrets Manager or SSM Parameter Store, not
          environment variables, for sensitive values.
    \item \textbf{Image hygiene} --- Multi-stage builds; strip tests and docs;
          \texttt{pip install --no-cache-dir}; remove \texttt{\_\_pycache\_\_}.
          Target $<$ 2 GB images when possible.
\end{enumerate}

% =========================================================
\section{Troubleshooting Checklist}
% =========================================================

\begin{longtable}{p{5.5cm} p{9.2cm}}
\toprule
\textbf{Symptom} & \textbf{Likely Cause / Fix} \\
\midrule
\texttt{exec format error} at cold start &
Image built for wrong architecture. Rebuild with
\texttt{--platform linux/amd64}. \\
\addlinespace
\texttt{Runtime.InvalidEntrypoint} &
\texttt{CMD} in Dockerfile does not match handler.
Must be \texttt{["file.function"]} (no \texttt{.py}). \\
\addlinespace
Cold start $>$ 10 s &
Image too large / model too large. Increase memory, use Provisioned
Concurrency, externalize weights to EFS, or switch to SageMaker. \\
\addlinespace
\texttt{Task timed out after 3.00 seconds} &
Default timeout is 3 s. Raise with \texttt{--timeout} (max 900). \\
\addlinespace
\texttt{[Errno 28] No space left on device} &
Writing to \texttt{/tmp} which defaults to 512 MB. Increase ephemeral
storage or write elsewhere. \\
\addlinespace
\texttt{AccessDeniedException} pulling image &
Lambda execution role needs ECR permissions, \emph{or} the ECR repo
policy must allow the Lambda service principal. \\
\addlinespace
Memory climbs invocation-over-invocation &
Module-level state accumulating. Clear caches in the handler or
recycle containers by bumping a no-op env var. \\
\bottomrule
\end{longtable}

% =========================================================
\section{Summary}
% =========================================================

The canonical Lambda + ECR deployment flow is:

\begin{enumerate}[leftmargin=*,nosep]
    \item Write handler with module-level model load.
    \item Write Dockerfile using \texttt{public.ecr.aws/lambda/python}.
    \item Create ECR repository.
    \item Build (\texttt{--platform linux/amd64}), tag, push.
    \item Create IAM execution role.
    \item Create Lambda function with \texttt{--package-type Image}.
    \item Expose via Function URL or API Gateway.
    \item Tune memory, timeout, ephemeral storage.
    \item Attach Provisioned Concurrency for latency-sensitive paths.
    \item Monitor with CloudWatch, X-Ray, Lambda Insights.
\end{enumerate}

Lambda + ECR is the right answer when your traffic is spiky, your model is
small-to-medium, and you want to pay nothing when idle --- which describes a
surprising fraction of real-world ML workloads.

\end{document}
